% from ega0-18-fr.tex
\oldpage[0\textsubscript{IV-1}]{61}

$w(zv(b'))=w(z)b'$ quels que soient $z\in L$ et $b'\in B'$ (en d'autres
termes, $w$ est un homomorphisme de $B'$-bimodules lorsqu'on munit $L$ de la
structure de $B'$-bimodule définie par $v$). La composition des morphismes est
définie par
$(u',v',w')\circ(u,v,w)=(u\circ u',v\circ v',w'\circ w)$, qui se justifie
aussitôt. Nous nous proposons de montrer que
\begin{equation}
\label{0.18.3.5.1-fr}
\tag{18.3.5.1}
  (A,B,L)\longmapsto\operatorname{Exan}_A(B,L)
\end{equation}
est un \emph{foncteur covariant} de la catégorie $\mathbf K$ dans la catégorie
$\mathbf{Ab}$ des groupes commutatifs. Il s'agit donc pour tout triplet
$(u,v,w)$ comme ci-dessus, de définir un homomorphisme de groupes commutatifs
\[
  (u,v,w)_*:\operatorname{Exan}_A(B,L)
  \longrightarrow\operatorname{Exan}_{A'}(B',L').
\]
En vertu de la définition des morphismes dans $\mathbf K$, on peut écrire
\[
  (u,v,w)=(1_{A'},1_{B'},w)\circ(1_{A'},v,1_L)
  \circ(u,1_B,1_L)
\]
où, dans le premier facteur, $L$ est muni de sa structure de $B'$-bimodule
définie par $v$~; nous définirons donc d'abord $(u,v,w)_*$ lorsque deux des
homomorphismes $u$, $v$, $w$ sont réduits à l'identité.
\end{env}

\begin{env}[18.3.6]
\label{0.18.3.6-fr}
Nous prendrons d'abord pour $(1_A,1_B,w)_*$ l'application
\begin{equation}
\label{0.18.3.6.1-fr}
\tag{18.3.6.1}
  w_*:\operatorname{Exan}_A(B,L)\longrightarrow
  \operatorname{Exan}_A(B,L')
\end{equation}
notée $T(w)$ dans (18.3.2)~; il est immédiat de vérifier que c'est un
homomorphisme de groupes, cette propriété s'exprimant par la commutativité de
diagrammes, transformés par $T$ de diagrammes analogues pour $L$ et $L'$.

L'application $(1_A,v,1_L)_*$ est l'application
\begin{equation}
\label{0.18.3.6.2-fr}
\tag{18.3.6.2}
  v^*:\operatorname{Exan}_A(B,L)\longrightarrow
  \operatorname{Exan}_A(B',L)
\end{equation}
définie de la façon suivante~: si $E$ est une $A$-extension de $B$ par $L$,
on a vu que $E\times_BB'$ est une $A$-extension de $B'$ par $L$ (18.2.5) et
que si l'on remplace $E$ par une $A$-extension $A$-équivalente $E'$,
$E'\times_BB'$ est $A$-équivalente à $E\times_BB'$~; l'image par $v^*$ de la
classe de $E$ est la classe de $E\times_BB'$. On vérifie aussitôt que si
$w:L\to L'$ est un homomorphisme de $B$-bimodules, le diagramme
\begin{equation}
\label{0.18.3.6.3-fr}
\tag{18.3.6.3}
  \xymatrix@C=3pc@R=2.4pc{
    \operatorname{Exan}_A(B,L)\ar[r]^{v^*}\ar[d]_{w_*}
      & \operatorname{Exan}_A(B',L)\ar[d]^{w_*}\\
    \operatorname{Exan}_A(B,L')\ar[r]_{v^*}
      & \operatorname{Exan}_A(B',L')
  }
\end{equation}
est commutatif, $L$ et $L'$ étant considérés comme $B'$-bimodules au moyen
de $v$ dans la colonne de droite. Remplaçant $L$ et $L'$ respectivement par
$L\times L$ et $L$, et $w$ par l'addition $s$ dans $L$, on en conclut que
$v^*$ est bien un \emph{homomorphisme de groupes}.
\end{env}

